% ============================================================================
% CHAPTER 3 — RESULTS
% ============================================================================
\chapter{Results}
\label{chap:results}

\section{Cohort / Dataset Characteristics}
\label{sec:cohort}

\placeholder{Present descriptive statistics and dataset characteristics.}

\section{Feature Results}
\label{sec:feature-results}

\placeholder{Present extracted features and important observations.}

% Example figure. Replace the filename with your own image.
% \begin{figure}[htbp]
%     \centering
%     \includegraphics[width=0.82\textwidth]{03-results/example_figure}
%     \caption{Descriptive caption explaining what is shown and, where relevant,
%     the experimental conditions.}
%     \label{fig:example}
% \end{figure}

\section{Model Performance}
\label{sec:model-performance}

\placeholder{Report performance metrics with uncertainty and validation details.}

% Example table.
\begin{table}[htbp]
    \centering
    \caption{Example model-performance table. Replace all values before submission.}
    \label{tab:example-performance}
    \begin{tabular}{lccc}
        \toprule
        Model & Accuracy & Precision & Recall \\
        \midrule
        Baseline & 0.00 & 0.00 & 0.00 \\
        Model A & 0.00 & 0.00 & 0.00 \\
        Model B & 0.00 & 0.00 & 0.00 \\
        \bottomrule
    \end{tabular}
\end{table}

As shown in \cref{tab:example-performance}, this is only a formatting example.
